
\documentclass[10pt,a4paper,twoside]{report}
\usepackage{fontspec}
\setmainfont{DejaVu Serif}[Scale=0.92]
\setsansfont{DejaVu Sans}[Scale=0.90]
\setmonofont{DejaVu Sans Mono}[Scale=0.80]
\usepackage[a4paper,top=22mm,bottom=24mm,left=24mm,right=20mm,headheight=14pt]{geometry}
\usepackage{xcolor}
\usepackage{graphicx}
\usepackage{longtable,tabularx,booktabs,array}
\usepackage{listings}
\usepackage{tikz}
\usetikzlibrary{shapes.geometric,positioning,calc,fit,backgrounds}
\usepackage{fancyhdr}
\usepackage{needspace}
\usepackage{enumitem}
\usepackage[hidelinks]{hyperref}
\usepackage{titlesec}
\usepackage{multicol}
\usepackage{rotating}
\sloppy\emergencystretch=3em

\definecolor{kw}{RGB}{20,70,160}
\definecolor{cm}{RGB}{90,110,90}
\definecolor{st}{RGB}{150,60,20}
\definecolor{rule}{RGB}{200,205,212}
\definecolor{okc}{RGB}{20,110,60}
\definecolor{badc}{RGB}{170,30,30}
\definecolor{panel}{RGB}{246,248,251}

\lstdefinestyle{js}{
  language=Java, basicstyle=\ttfamily\scriptsize, keywordstyle=\color{kw},
  commentstyle=\color{cm}\itshape, stringstyle=\color{st},
  breaklines=true, breakatwhitespace=false, columns=fullflexible, keepspaces=true,
  showstringspaces=false, tabsize=2, numbers=left, numberstyle=\ttfamily\tiny\color{gray},
  numbersep=7pt, xleftmargin=20pt, framexleftmargin=16pt, frame=single, rulecolor=\color{rule},
  backgroundcolor=\color{panel}, aboveskip=7pt, belowskip=7pt, upquote=true,
  literate={·}{{\textperiodcentered}}1 {—}{{---}}1 {–}{{--}}1 {…}{{\ldots}}1 {≤}{{$\leq$}}1
           {≥}{{$\geq$}}1 {×}{{$\times$}}1 {→}{{$\rightarrow$}}1 {²}{{$^2$}}1 {°}{{$^\circ$}}1
           {₀}{{$_0$}}1 {"}{{"}}1 {"}{{"}}1 {'}{{'}}1 {'}{{'}}1
}
\lstdefinestyle{jstiny}{style=js, basicstyle=\ttfamily\tiny, numbers=none}
\lstdefinestyle{plain}{style=js, keywordstyle=\color{black}, numbers=none, language={}}

\pagestyle{fancy}\fancyhf{}
\renewcommand{\headrulewidth}{0.3pt}\renewcommand{\footrulewidth}{0.3pt}
\fancyhead[LE,RO]{\sffamily\scriptsize RDML reader — validation report}
\fancyhead[RE,LO]{\sffamily\scriptsize\leftmark}
\fancyfoot[C]{\sffamily\scriptsize Page \thepage\ of \pageref{LastPage}}
\usepackage{lastpage}
\fancypagestyle{plain}{\fancyhf{}\fancyfoot[C]{\sffamily\scriptsize Page \thepage\ of \pageref{LastPage}}
  \renewcommand{\headrulewidth}{0pt}\renewcommand{\footrulewidth}{0.3pt}}

\titleformat{\chapter}[display]{\sffamily\LARGE\bfseries}{\sffamily\normalsize\MakeUppercase{\chaptertitlename\ \thechapter}}{6pt}{}
\titlespacing*{\chapter}{0pt}{-14pt}{18pt}
\titleformat{\section}{\sffamily\large\bfseries}{\thesection}{0.7em}{}
\titleformat{\subsection}{\sffamily\normalsize\bfseries}{\thesubsection}{0.7em}{}
\setlength{\parindent}{0pt}\setlength{\parskip}{4pt}
\setlength{\tabcolsep}{4pt}\renewcommand{\arraystretch}{1.16}
\newcommand{\PASS}{\textcolor{okc}{\textbf{pass}}}
\newcommand{\FAILED}{\textcolor{badc}{\textbf{fail}}}
\newcommand{\fn}[1]{\texttt{#1}}
\newcommand{\AThreeStart}{\clearpage\pdfpagewidth=420mm\pdfpageheight=297mm
  \newgeometry{left=18mm,right=18mm,top=22mm,bottom=22mm,headheight=14pt,headsep=8mm,footskip=14mm}%
  \setlength{\textwidth}{384mm}\setlength{\textheight}{230mm}%
  \setlength{\oddsidemargin}{\dimexpr 18mm-1in\relax}%
  \setlength{\evensidemargin}{\dimexpr 18mm-1in\relax}%
  \setlength{\topmargin}{\dimexpr 22mm-1in\relax}%
  \setlength{\headheight}{14pt}\setlength{\headsep}{8mm}\setlength{\footskip}{14mm}%
  \setlength{\marginparwidth}{0pt}\setlength{\linewidth}{\textwidth}%
  \setlength{\columnwidth}{\textwidth}\hsize=\textwidth
  \setlength{\headwidth}{\textwidth}}
\newcommand{\AThreeEnd}{\clearpage\pdfpagewidth=210mm\pdfpageheight=297mm\restoregeometry
  \setlength{\headwidth}{\textwidth}}
\tikzset{
  box/.style={draw=rule!60!black, fill=panel, rounded corners=2pt, align=center,
              font=\sffamily\footnotesize, inner sep=5pt, minimum height=9mm, text width=34mm},
  wide/.style={box, text width=52mm},
  narrow/.style={box, text width=25mm},
  dec/.style={draw=rule!60!black, fill=white, diamond, aspect=2.4, align=center,
              font=\sffamily\scriptsize, inner sep=1pt, text width=30mm},
  io/.style={box, fill=white, draw=kw!70},
  ok/.style={box, fill=okc!8, draw=okc!60},
  bad/.style={box, fill=badc!8, draw=badc!60},
  fl/.style={->, draw=rule!50!black, thick},
  lbl/.style={font=\sffamily\scriptsize, inner sep=1.5pt, fill=white, align=center}
}

\begin{document}

\begin{titlepage}\centering\sffamily
\vspace*{18mm}
{\Huge\bfseries RDML reader\par}
\vspace{4mm}
{\LARGE Function-by-function validation report\par}
\vspace{10mm}
{\large qPCR raw values, open formats, quality control and forensics\par}
\vspace{2mm}
{\large LightCycler 480 \textperiodcentered{} QuantStudio 3/5 \textperiodcentered{} RDML 1.0--1.4\par}
\vspace{16mm}
\begin{minipage}{0.86\textwidth}\centering\small
\begin{tabular}{@{}ll@{}}
\toprule
Build under test & \texttt{qpcr\_qc\_forensics.html} \\
SHA-256 & \texttt{6e14cc49153ad0662b94c47218f213051c5128fec78cc8a479ce00f17f391f4e} \\
Size & 860,756 bytes \\
Symbols validated & 15 functions, 2 constants \\
Test records & 89 \\
Corpus & 8 instrument and exchange files \\
Browser faults during capture & 0 \\
Capture run & \texttt{2026-09-24T08:44:56.803Z} \\
\bottomrule
\end{tabular}
\end{minipage}
\vfill
{\small This report states, for every function of the RDML reader, what it was given, where that
input came from, what it returned, and how the verdict was reached. Inputs that were constructed
for the test are marked as such and printed in full.\par}
\vspace{8mm}
\end{titlepage}
\tableofcontents
\clearpage

\chapter{Purpose, scope and reading guide}

\section{What this document is}
This is the validation record of the RDML reader added to the page on 24 September 2026. It is
written to be checked rather than believed: every number in it was produced by a harness that is
printed in Appendix~\ref{app:harness}, run against the build whose SHA-256 is on the title page,
and the harness writes a machine-readable record (\texttt{validation\_data.json}) from which this
document is generated. Re-running the harness regenerates every table here.

\section{Why the provenance of the input matters}
The functions under test were written together with most of their test inputs. That is a conflict
of interest, and the only honest response is to declare, for every single input, where it came
from. Three classes are used throughout, and every input table names its class:

\begin{longtable}{@{}p{10mm}p{42mm}p{\dimexpr\linewidth-58mm-16pt\relax}@{}}
\toprule
\textbf{Class} & \textbf{Meaning} & \textbf{What it implies for the evidence} \\
\midrule\endhead

\textbf{R} & Real file, unmodified & A file supplied by the laboratory or by the RDML workshop, used byte for byte. Its SHA-256 is in Table 2. \\
\addlinespace

\textbf{D} & Derived from a real file & An object produced from a real file by a stated operation of the code under test, e.g. the first \texttt{<run>} element of a parsed document. \\
\addlinespace

\textbf{S} & Synthetic, written for the test & A literal value constructed for this report. Every synthetic value is printed in full so the reader can reproduce it. \\
\addlinespace

\bottomrule
\end{longtable}

Class~R inputs carry the whole weight of the validation: they were produced by instruments and by
the RDML-Tools of the Freising workshop, long before this reader existed, and neither their content
nor their structure could be arranged to suit it. Class~D inputs are real data that the code under
test has already transformed, so they test a later stage against the output of an earlier one.
Class~S inputs exist only for boundary conditions --- an empty string, a malformed document, a null
argument --- where no real file offers the case. They are printed verbatim so that nothing is hidden
in a description.

\section{The strongest single piece of evidence}
One experiment in the corpus, \texttt{CORPUS-RDML-1}, exists both as the instrument's own
container (\texttt{.eds}) and as an RDML export of the same run. The page already decodes the first
by a path that has nothing to do with the reader under test. Reading both and comparing the Cq of
every reaction is therefore an external check that no amount of self-written test data could
provide. It is reported in section~\ref{sec:oracle}.

\section{How to read a function section}
Each function gets: what it is for; its signature and its place in the call graph; its complete
source as it stands in the build under test, with the production line numbers; a table of every
input it was given, with the provenance class; a table of what it returned; a flow diagram where
the control flow justifies one; and the verdict. Nothing is abbreviated. Where a returned value is
large it is printed in full in an appendix rather than truncated in place.

\chapter{The corpus and how it was obtained}\label{ch:corpus}

\section{Files}
Eight files, none of them produced by this project. Five are the exercise files of the RDML-Tools
workshop \emph{Web-based analysis of qPCR data} (Freising, March 2023); two are exchange exports
supplied by the laboratory; one is the instrument container that pairs with the first of those.

\begin{longtable}{@{}p{56mm}r p{\dimexpr\linewidth-56mm-20mm-28pt\relax}@{}}
\toprule
\textbf{File} & \textbf{Bytes} & \textbf{SHA-256} \\
\midrule\endhead

\texttt{2026-\allowbreak{}06-\allowbreak{}15\_\allowbreak{}RAS\_\allowbreak{}FFPE.\allowbreak{}eds} & 1,756,735 & {\ttfamily\scriptsize d2750711198559b16cc963d9f9c074c1\newline 4b8a248c3fbd14436852884870d6be6c} \\

\multicolumn{3}{@{}p{\linewidth}@{}}{\small Instrument container, QuantStudio. Decoded by the page's existing \texttt{.eds} path, not by the reader under test. Reference for the cross-format oracle.}\\
\addlinespace

\texttt{2026-\allowbreak{}06-\allowbreak{}15\_\allowbreak{}RAS\_\allowbreak{}FFPE.\allowbreak{}rdml} & 69,062 & {\ttfamily\scriptsize 688ef2eeec2e0fd350e8d2e585933e86\newline 5b2eed7ae82e56bb0c6c0c4f60ddc384} \\

\multicolumn{3}{@{}p{\linewidth}@{}}{\small RDML 1.2 export of the run above. Names an instrument; carries Cq; no re-analysis tool.}\\
\addlinespace

\texttt{Gene\allowbreak{}Expression\_\allowbreak{}dd\allowbreak{}Ct\_\allowbreak{}Fast\_\allowbreak{}Adv\_\allowbreak{}MMx\_\allowbreak{}10u\allowbreak{}L.\allowbreak{}rdml} & 23,248 & {\ttfamily\scriptsize a26b7c8a7a596d3376d776ff32990eb8\newline 0b855d5e657d274fbe56c475085a5c30} \\

\multicolumn{3}{@{}p{\linewidth}@{}}{\small RDML 1.2 export of a QuantStudio demonstration run.}\\
\addlinespace

\texttt{example\_\allowbreak{}1\_\allowbreak{}raw.\allowbreak{}rdml} & 252,382 & {\ttfamily\scriptsize a29e4918383b5199a5757047d884c149\newline f21650a352bb9c487ccc398769b37a3a} \\

\multicolumn{3}{@{}p{\linewidth}@{}}{\small Workshop file, RDML 1.3. Amplification and melting fluorescence only; no Cq anywhere.}\\
\addlinespace

\texttt{example\_\allowbreak{}2\_\allowbreak{}tm\_\allowbreak{}annotated.\allowbreak{}rdml} & 252,547 & {\ttfamily\scriptsize e018dbad7ec4071af882543929dad171\newline e7297517e4e0f6a39c8cae4d8f9ba6d7} \\

\multicolumn{3}{@{}p{\linewidth}@{}}{\small Workshop file. The same data with melting temperatures annotated on the targets.}\\
\addlinespace

\texttt{example\_\allowbreak{}3\_\allowbreak{}linregpcr.\allowbreak{}rdml} & 260,964 & {\ttfamily\scriptsize 4806d1466a741c7f31141ce2d5cf0a02\newline da5d329180e392d0277ddc12c82c8845} \\

\multicolumn{3}{@{}p{\linewidth}@{}}{\small Workshop file. The same data after LinRegPCR: Cq, N$_0$, efficiency, exclusions and notes.}\\
\addlinespace

\texttt{example\_\allowbreak{}4\_\allowbreak{}interrun.\allowbreak{}rdml} & 262,745 & {\ttfamily\scriptsize 7545883ce03d305207f8b73d40b403a2\newline a919f203908e8f4bbe27ae8b3019d4b2} \\

\multicolumn{3}{@{}p{\linewidth}@{}}{\small Workshop file. The same data again after between-plate correction.}\\
\addlinespace

\texttt{example\_\allowbreak{}LCGreen\_\allowbreak{}raw.\allowbreak{}rdml} & 52,635 & {\ttfamily\scriptsize 73f238867ae4dcdd12511b9f3519fdd5\newline f8e6653eac3a97cbae451da4bd7e9293} \\

\multicolumn{3}{@{}p{\linewidth}@{}}{\small Workshop file. Single dye, one run, deliberate artefacts; no Cq.}\\
\addlinespace

\bottomrule
\end{longtable}

\section{What the corpus contains, as read}
The table below is the output of \fn{rdmlParse} over the whole corpus and is repeated in
section~\ref{sec:rdmlParse} with the per-run detail. It is placed here because it also serves as the
description of the corpus.

\begin{longtable}{@{}p{40mm}c c r r r r p{22mm}@{}}
\toprule
\textbf{File} & \textbf{ver} & \textbf{runs} & \textbf{wells} & \textbf{curves} & \textbf{melt} & \textbf{Cq} & \textbf{state} \\
\midrule\endhead

\texttt{2026-\allowbreak{}06-\allowbreak{}15\_\allowbreak{}RAS\_\allowbreak{}FFPE} & 1.2 & 1 & 192 & 96 & 0 & 34 & vendor-export \\

\texttt{Gene\allowbreak{}Expression\_\allowbreak{}dd\allowbreak{}Ct\_\allowbreak{}Fast\_\allowbreak{}Adv\_\allowbreak{}MMx\_\allowbreak{}10u\allowbreak{}L} & 1.2 & 1 & 72 & 72 & 0 & 66 & -- \\

\texttt{example\_\allowbreak{}1\_\allowbreak{}raw} & 1.3 & 3 & 285 & 285 & 452 & 0 & raw \\

\texttt{example\_\allowbreak{}2\_\allowbreak{}tm\_\allowbreak{}annotated} & 1.3 & 3 & 285 & 285 & 452 & 0 & -- \\

\texttt{example\_\allowbreak{}3\_\allowbreak{}linregpcr} & 1.3 & 3 & 285 & 285 & 452 & 266 & third-party \\

\texttt{example\_\allowbreak{}4\_\allowbreak{}interrun} & 1.3 & 3 & 285 & 285 & 452 & 266 & -- \\

\texttt{example\_\allowbreak{}LCGreen\_\allowbreak{}raw} & 1.3 & 1 & 50 & 50 & 152 & 0 & -- \\

\bottomrule
\end{longtable}
\small ``melt'' is the number of melting acquisition points of the first run of the file; ``Cq'' is
the number of reactions carrying a crossing point after the ceiling rule of
section~\ref{sec:ceiling} has been applied.\normalsize
